% Краткое сообщение Марине: Марина, основное внимание стоит уделить завершению решения до требуемого ответа, выбору формулы именно под вопрос задачи и аккуратной работе со знаками. В задании про коэффициенты параболы вычисления выполнены верно, но важно обязательно записывать итоговую формулу, если она требуется условием.
\documentclass[12pt,a4paper]{article}

\usepackage[margin=2cm]{geometry}
\usepackage[T2A]{fontenc}
\usepackage[utf8]{inputenc}
\usepackage[russian]{babel}
\usepackage{amsmath,amssymb,mathtools}
\usepackage{array,booktabs,longtable}
\usepackage{xcolor}
\usepackage{hyperref}
\usepackage{tikz}

\hypersetup{
  colorlinks=true,
  linkcolor=blue!55!black,
  urlcolor=blue!55!black
}

\setlength{\parindent}{0pt}
\setlength{\parskip}{0.55em}
\renewcommand{\arraystretch}{1.25}

\begin{document}
\hypersetup{pageanchor=false}

% -------------------- Титульный лист --------------------
\begin{titlepage}
\centering
\vspace*{0.24\textheight}
{\LARGE\bfseries Ревью домашней работы\par}
\vspace{1.7cm}
{\large Автор: Лёвин Артём Александрович\par}
\vspace{0.8cm}
{\large 16 сентября 2026 года\par}
\vfill
\begin{tikzpicture}[x=1cm,y=1cm]
  \draw[blue!45!black, line width=0.7pt]
    (0,0) .. controls (2,0.55) and (4,-0.55) .. (6,0);
\end{tikzpicture}
\end{titlepage}
\hypersetup{pageanchor=true}
\pagenumbering{arabic}

% -------------------- Сводная таблица --------------------
\newpage
\section*{Сводная таблица}
\small
\setlength{\tabcolsep}{3pt}
\begin{longtable}{@{}>{\centering\arraybackslash}p{0.07\linewidth}p{0.17\linewidth}p{0.26\linewidth}p{0.17\linewidth}p{0.20\linewidth}@{}}
\toprule
\textbf{№} & \textbf{Тема} & \textbf{Суть ошибки} & \textbf{Ваш ответ} & \textbf{Правильный ответ} \\
\midrule
\endfirsthead
\toprule
\textbf{№} & \textbf{Тема} & \textbf{Суть ошибки} & \textbf{Ваш ответ} & \textbf{Правильный ответ} \\
\midrule
\endhead
\hyperref[task:3]{3} & Пересечение двух прямых & Решение не доведено до системы и координат точки пересечения. & Не указан. & $(3;5)$. \\
\addlinespace
\hyperref[task:5]{5} & Вершина параболы & Найден дискриминант, но он не даёт абсциссу вершины; требуемый ответ не получен. & Не указан. & $x_{\text{в}}=2$, вершина справа от оси $Oy$. \\
\addlinespace
\hyperref[task:6]{6} & Пересечение прямой и параболы & Решение не выполнено: не составлено уравнение равенства ординат. & Не указан. & $(0;1)$ и $(4;5)$. \\
\addlinespace
\hyperref[task:7]{7} & Параметр в уравнении прямой & При переносе единицы допущена ошибка со знаком. & $m=-3$. & $m=3$. \\
\addlinespace
\hyperref[task:8]{8} & Горизонтальная прямая и парабола & Решение не выполнено; не использовано условие ровно одной общей точки. & Не указан. & $m=-1$. \\
\addlinespace
\hyperref[task:10]{10} & Коэффициенты параболы & Коэффициенты найдены верно, но обязательная итоговая формула параболы не записана. & $b=-12$, $c=-2$; формула не записана. & $y=2x^2-12x-2$. \\
\bottomrule
\end{longtable}
\normalsize

% -------------------- Разделитель --------------------
\newpage
\thispagestyle{empty}
\vspace*{\fill}
\begin{center}
{\LARGE\bfseries Разбор ошибок}
\end{center}
\vspace*{\fill}
\newpage

% ==================== Задание 3 ====================
\phantomsection\label{task:3}
\section*{Задание №3}

\textbf{Текст задания.}
Найдите точку пересечения прямых
\[
y=3x-4, \qquad y=-x+8.
\]

\textbf{Ваше решение.}
В тетради записаны обе прямые и начат эскиз их пересечения. Числовое решение системы и координаты точки пересечения не записаны.

\textbf{Ваш ответ.}
Не указан.

\textcolor{red}{\textbf{Ошибка:}} решение не доведено до координат точки пересечения.

\textbf{Где именно возникает ошибка.}

\emph{Последний корректный шаг:} обе прямые записаны верно.

\emph{Первый неполный шаг:} после записи уравнений не выполнено приравнивание их правых частей. Эскиз без подписанного масштаба не позволяет надёжно получить точные координаты точки пересечения.

\textbf{Почему это неверно.}
Точка пересечения должна одновременно лежать на обеих прямых. Поэтому при одном и том же значении $x$ её ордината $y$ должна быть одинаковой в обоих уравнениях. Значит, необходимо приравнять выражения $3x-4$ и $-x+8$, найти $x$, а затем подставить найденное значение хотя бы в одно из исходных уравнений. Пока координаты не вычислены и не записаны, требование задачи не выполнено.

\textcolor{blue!60!black}{\textbf{Ключевая идея:}} в точке пересечения значения $y$ у двух прямых равны.

\textbf{Исправление от последнего правильного шага.}
Из записанных уравнений получаем
\[
3x-4=-x+8.
\]
Перенесём $-x$ в левую часть, а $-4$ --- в правую:
\[
3x+x=8+4.
\]
Отсюда
\[
4x=12,
\qquad
x=3.
\]
Теперь найдём ординату:
\[
y=3\cdot 3-4=9-4=5.
\]

\textbf{Полное правильное решение.}
Точка пересечения принадлежит обеим прямым, поэтому
\[
3x-4=-x+8.
\]
Решаем линейное уравнение:
\[
3x+x=8+4,
\]
\[
4x=12,
\]
\[
x=3.
\]
Подставляем $x=3$ в первое уравнение:
\[
y=3\cdot 3-4=5.
\]
Следовательно, точка пересечения имеет координаты
\[
(3;5).
\]

\textbf{Проверка.}
Подставим $x=3$ во вторую прямую:
\[
y=-3+8=5.
\]
Обе прямые при $x=3$ дают одно и то же значение $y=5$, значит, точка $(3;5)$ действительно лежит на обеих прямых.

\textcolor{teal!60!black}{\textbf{Итог:}} точка пересечения прямых --- $(3;5)$.

\textbf{Задача на закрепление.}
Найдите точку пересечения прямых
\[
y=2x+1, \qquad y=-x+7.
\]

% ==================== Задание 5 ====================
\newpage
\phantomsection\label{task:5}
\section*{Задание №5}

\textbf{Текст задания.}
Для параболы
\[
y=2x^2-8x+3
\]
найдите абсциссу вершины и укажите, справа или слева от оси $Oy$ она находится.

\textbf{Ваше решение.}
Записано уравнение параболы и вычислен дискриминант:
\[
D=64-24=40.
\]
Дальнейшего решения и ответа нет.

\textbf{Ваш ответ.}
Не указан.

\textcolor{red}{\textbf{Ошибка:}} выбран вычислительный шаг, который не отвечает на поставленный вопрос, и решение не завершено.

\textbf{Где именно возникает ошибка.}

\emph{Последний корректный шаг:} дискриминант вычислен арифметически верно: $D=40$.

\emph{Первый неполный шаг:} после вычисления дискриминанта решение остановлено, хотя нужно было найти абсциссу вершины по коэффициентам $a$ и $b$.

\textbf{Почему это неверно.}
Дискриминант квадратного трёхчлена используется прежде всего для исследования и нахождения его нулей. В задаче спрашивается не о корнях, а о вершине параболы. Для параболы $y=ax^2+bx+c$ абсцисса вершины находится по формуле
\[
x_{\text{в}}=-\frac{b}{2a}.
\]
После этого знак найденной абсциссы сразу определяет положение вершины относительно оси $Oy$: положительная абсцисса означает положение справа, отрицательная --- слева.

\textcolor{blue!60!black}{\textbf{Ключевая идея:}} для вершины параболы нужно использовать формулу $x_{\text{в}}=-\dfrac{b}{2a}$, а не дискриминант.

\textbf{Исправление от последнего правильного шага.}
В данном уравнении
\[
a=2, \qquad b=-8.
\]
Поэтому
\[
x_{\text{в}}=-\frac{-8}{2\cdot 2}=\frac{8}{4}=2.
\]
Так как $2>0$, вершина расположена справа от оси $Oy$.

\textbf{Полное правильное решение.}
Для параболы
\[
y=2x^2-8x+3
\]
коэффициенты равны
\[
a=2, \qquad b=-8, \qquad c=3.
\]
Абсцисса вершины:
\[
x_{\text{в}}=-\frac{b}{2a}.
\]
Подставляем $a=2$ и $b=-8$:
\[
x_{\text{в}}=-\frac{-8}{2\cdot 2}=2.
\]
Абсцисса положительна, следовательно, вершина находится справа от оси $Oy$.

\textbf{Проверка.}
Выделим полный квадрат:
\[
y=2(x^2-4x)+3
=2\bigl((x-2)^2-4\bigr)+3
=2(x-2)^2-5.
\]
Из этой формы непосредственно видно, что вершина имеет абсциссу $2$. Это подтверждает найденный результат.

\textcolor{teal!60!black}{\textbf{Итог:}} $x_{\text{в}}=2$; вершина расположена справа от оси $Oy$.

\textbf{Задача на закрепление.}
Для параболы
\[
y=3x^2+12x-5
\]
найдите абсциссу вершины и укажите, справа или слева от оси $Oy$ она находится.

% ==================== Задание 6 ====================
\newpage
\phantomsection\label{task:6}
\section*{Задание №6}

\textbf{Текст задания.}
Найдите точки пересечения прямой и параболы:
\[
y=x+1, \qquad y=x^2-3x+1.
\]

\textbf{Ваше решение.}
Рядом с номером задания стоит знак вопроса; вычислений нет.

\textbf{Ваш ответ.}
Не указан.

\textcolor{red}{\textbf{Ошибка:}} решение задания не выполнено.

\textbf{Где именно возникает ошибка.}

\emph{Последний корректный шаг:} выполненных математических шагов нет; условие только отмечено.

\emph{Первый неполный шаг:} не составлено уравнение, выражающее равенство ординат прямой и параболы в точках пересечения.

\textbf{Почему это неверно.}
Если точка лежит одновременно на прямой и на параболе, то при одном и том же $x$ обе формулы должны давать одинаковое значение $y$. Поэтому необходимо приравнять правые части двух уравнений. Полученное квадратное уравнение даст все возможные абсциссы точек пересечения. После этого для каждой абсциссы нужно найти соответствующую ординату.

\textcolor{blue!60!black}{\textbf{Ключевая идея:}} точки пересечения находятся из равенства $x+1=x^2-3x+1$.

\textbf{Исправление от последнего правильного шага.}
Приравниваем правые части:
\[
x+1=x^2-3x+1.
\]
Переносим всё в одну часть:
\[
0=x^2-4x.
\]
Выносим $x$ за скобку:
\[
x(x-4)=0.
\]
Отсюда
\[
x=0 \quad \text{или} \quad x=4.
\]
Находим соответствующие значения $y$ по формуле прямой $y=x+1$:
\[
x=0 \Rightarrow y=1,
\]
\[
x=4 \Rightarrow y=5.
\]

\textbf{Полное правильное решение.}
В точке пересечения значения $y$ совпадают, поэтому
\[
x+1=x^2-3x+1.
\]
Вычитаем $x+1$ из обеих частей:
\[
x^2-3x+1-x-1=0,
\]
\[
x^2-4x=0.
\]
Разлагаем на множители:
\[
x(x-4)=0.
\]
Произведение равно нулю, если хотя бы один множитель равен нулю:
\[
x=0 \quad \text{или} \quad x=4.
\]
Для $x=0$:
\[
y=0+1=1.
\]
Получаем точку $(0;1)$.

Для $x=4$:
\[
y=4+1=5.
\]
Получаем точку $(4;5)$.

\textbf{Проверка.}
Для точки $(0;1)$ по параболе:
\[
0^2-3\cdot 0+1=1.
\]
Для точки $(4;5)$ по параболе:
\[
4^2-3\cdot 4+1=16-12+1=5.
\]
Обе точки удовлетворяют и прямой, и параболе.

\textcolor{teal!60!black}{\textbf{Итог:}} точки пересечения --- $(0;1)$ и $(4;5)$.

\textbf{Задача на закрепление.}
Найдите точки пересечения прямой и параболы:
\[
y=2x-1, \qquad y=x^2-x-1.
\]

% ==================== Задание 7 ====================
\newpage
\phantomsection\label{task:7}
\section*{Задание №7}

\textbf{Текст задания.}
При каком значении параметра $m$ точка $(2;7)$ лежит на прямой
\[
y=mx+1?
\]

\textbf{Ваше решение.}
Записано:
\[
7=2m+1,
\]
затем
\[
2m=-7+1,
\]
после чего получено
\[
m=-3.
\]

\textbf{Ваш ответ.}
$m=-3$.

\textcolor{red}{\textbf{Ошибка:}} неверно изменён знак при преобразовании линейного уравнения.

\textbf{Где именно возникает ошибка.}

\emph{Последний корректный шаг:}
\[
7=2m+1.
\]
Подстановка координат точки выполнена верно.

\emph{Первый неправильный шаг:}
\[
2m=-7+1.
\]
Из равенства $7=2m+1$ это преобразование не следует.

\textbf{Почему этот шаг неверен.}
Чтобы оставить слагаемое $2m$ отдельно, нужно вычесть $1$ из обеих частей равенства:
\[
7-1=2m+1-1.
\]
Тогда получается
\[
6=2m.
\]
Можно записать то же самое как
\[
2m=7-1.
\]
Число $7$ не переносится через знак равенства: оно остаётся в той же части уравнения. Поэтому менять знак у $7$ оснований нет. Ошибка со знаком меняет $6$ на $-6$ и приводит к противоположному значению параметра.

\textcolor{blue!60!black}{\textbf{Ключевая идея:}} при преобразовании уравнения выполняйте одно и то же действие с обеими частями; здесь нужно вычесть $1$, а не менять знак у $7$.

\textbf{Исправление от последнего правильного шага.}
Начинаем с верного равенства
\[
7=2m+1.
\]
Вычитаем $1$ из обеих частей:
\[
7-1=2m.
\]
Получаем
\[
6=2m,
\]
следовательно,
\[
m=3.
\]

\textbf{Полное правильное решение.}
Если точка $(2;7)$ лежит на прямой $y=mx+1$, её координаты должны удовлетворять уравнению прямой. Подставляем $x=2$ и $y=7$:
\[
7=2m+1.
\]
Вычитаем $1$ из обеих частей:
\[
6=2m.
\]
Делим обе части на $2$:
\[
m=3.
\]

\textbf{Проверка.}
Подставим $m=3$ и $x=2$ в уравнение прямой:
\[
y=3\cdot 2+1=7.
\]
Получено именно значение ординаты данной точки, значит, параметр найден верно.

\textcolor{teal!60!black}{\textbf{Итог:}} $m=3$.

\textbf{Задача на закрепление.}
При каком значении параметра $m$ точка $(4;9)$ лежит на прямой
\[
y=mx+1?
\]

% ==================== Задание 8 ====================
\newpage
\phantomsection\label{task:8}
\section*{Задание №8}

\textbf{Текст задания.}
Для какого значения $m$ горизонтальная прямая
\[
y=m
\]
имеет с параболой
\[
y=(x-2)^2-1
\]
ровно одну общую точку?

\textbf{Ваше решение.}
Рядом с номером задания стоит знак вопроса; вычислений нет.

\textbf{Ваш ответ.}
Не указан.

\textcolor{red}{\textbf{Ошибка:}} не использовано условие о ровно одной общей точке прямой и параболы.

\textbf{Где именно возникает ошибка.}

\emph{Последний корректный шаг:} выполненных математических шагов нет; условие только отмечено.

\emph{Первый неполный шаг:} не составлено уравнение пересечения
\[
m=(x-2)^2-1.
\]

\textbf{Почему это неверно.}
Количество общих точек прямой и параболы равно количеству действительных решений уравнения, полученного приравниванием их значений $y$. После преобразования получаем
\[
(x-2)^2=m+1.
\]
Левая часть является квадратом и потому неотрицательна. Если $m+1>0$, уравнение имеет два решения: $x-2$ может быть положительным или отрицательным. Если $m+1<0$, действительных решений нет. Ровно одно решение возникает только тогда, когда квадрат равен нулю.

\textcolor{blue!60!black}{\textbf{Ключевая идея:}} горизонтальная прямая имеет с этой параболой ровно одну общую точку, когда проходит через её вершину.

\textbf{Исправление от последнего правильного шага.}
Приравниваем значения $y$:
\[
m=(x-2)^2-1.
\]
Переносим $-1$ в левую часть:
\[
(x-2)^2=m+1.
\]
Для ровно одного действительного значения $x$ необходимо
\[
m+1=0.
\]
Поэтому
\[
m=-1.
\]

\textbf{Полное правильное решение.}
Парабола
\[
y=(x-2)^2-1
\]
записана в вершинной форме. Её вершина имеет координаты $(2;-1)$. Прямая $y=m$ горизонтальна. Такая прямая пересекает параболу ровно один раз только в случае касания в вершине, то есть когда её высота равна ординате вершины:
\[
m=-1.
\]

То же самое можно получить алгебраически:
\[
m=(x-2)^2-1,
\]
\[
(x-2)^2=m+1.
\]
Ровно одно решение возможно при
\[
m+1=0,
\]
то есть
\[
m=-1.
\]

\textbf{Проверка.}
При $m=-1$ получаем
\[
-1=(x-2)^2-1.
\]
Тогда
\[
(x-2)^2=0,
\]
откуда единственное решение
\[
x=2.
\]
Следовательно, общая точка действительно одна: $(2;-1)$.

\textcolor{teal!60!black}{\textbf{Итог:}} $m=-1$.

\textbf{Задача на закрепление.}
Для какого значения $m$ горизонтальная прямая
\[
y=m
\]
имеет с параболой
\[
y=(x+3)^2+2
\]
ровно одну общую точку?

% ==================== Задание 10 ====================
\newpage
\phantomsection\label{task:10}
\section*{Задание №10}

\textbf{Текст задания.}
Парабола
\[
y=ax^2+bx+c
\]
проходит через точку $(0;-2)$, абсцисса её вершины равна $3$, а $a=2$. Найдите $b$ и $c$, затем запишите формулу параболы.

\textbf{Ваше решение.}
Верно найдено
\[
c=-2
\]
из условия прохождения через точку $(0;-2)$, а также
\[
b=-2ax_{\text{в}}=-2\cdot 2\cdot 3=-12.
\]
В конце записаны значения $a=2$, $b=-12$, $c=-2$, но итоговая формула параболы не записана.

\textbf{Ваш ответ.}
$b=-12$, $c=-2$; формула параболы не указана.

\textcolor{red}{\textbf{Ошибка:}} пропущен обязательный заключительный элемент ответа --- формула параболы.

\textbf{Где именно возникает ошибка.}

\emph{Последний корректный шаг:} коэффициенты найдены верно:
\[
a=2, \qquad b=-12, \qquad c=-2.
\]

\emph{Первый неполный шаг:} после нахождения коэффициентов они не подставлены в общую формулу $y=ax^2+bx+c$.

\textbf{Почему это неверно.}
Математические вычисления коэффициентов выполнены правильно, однако условие содержит два требования: сначала найти $b$ и $c$, затем записать конкретную формулу параболы. Перечень коэффициентов не заменяет вторую часть ответа. Чтобы завершить решение, нужно подставить найденные значения $a$, $b$, $c$ в общую формулу.

\textcolor{blue!60!black}{\textbf{Ключевая идея:}} после нахождения коэффициентов обязательно выполнить последнюю подстановку в $y=ax^2+bx+c$.

\textbf{Исправление от последнего правильного шага.}
Из уже найденных значений
\[
a=2, \qquad b=-12, \qquad c=-2
\]
получаем
\[
y=2x^2-12x-2.
\]

\textbf{Полное правильное решение.}
Так как парабола проходит через точку $(0;-2)$, подставляем $x=0$ и $y=-2$ в общую формулу:
\[
-2=a\cdot 0^2+b\cdot 0+c.
\]
Отсюда
\[
c=-2.
\]

Абсцисса вершины параболы находится по формуле
\[
x_{\text{в}}=-\frac{b}{2a}.
\]
По условию $x_{\text{в}}=3$ и $a=2$, поэтому
\[
3=-\frac{b}{2\cdot 2}=-\frac{b}{4}.
\]
Умножаем обе части на $4$:
\[
12=-b,
\]
следовательно,
\[
b=-12.
\]
Теперь подставляем
\[
a=2, \qquad b=-12, \qquad c=-2
\]
в формулу $y=ax^2+bx+c$:
\[
y=2x^2-12x-2.
\]

\textbf{Проверка.}
Проверим прохождение через заданную точку:
\[
y(0)=2\cdot 0^2-12\cdot 0-2=-2.
\]
Точка $(0;-2)$ принадлежит параболе.

Проверим абсциссу вершины:
\[
x_{\text{в}}=-\frac{-12}{2\cdot 2}=\frac{12}{4}=3.
\]
Оба условия выполняются.

\textcolor{teal!60!black}{\textbf{Итог:}} $b=-12$, $c=-2$, формула параболы
\[
y=2x^2-12x-2.
\]

\textbf{Задача на закрепление.}
Парабола
\[
y=ax^2+bx+c
\]
проходит через точку $(0;5)$, абсцисса её вершины равна $-2$, а $a=1$. Найдите $b$ и $c$, затем запишите формулу параболы.

\vfill
\begin{center}
\footnotesize Лёвин Артём Александрович эксклюзивно для Марины
\end{center}

\end{document}
